%%% Реализация библиографии через BibLaTeX (современный подход)

%%% Пакеты %%%
\usepackage[backend=biber,
            style=gost-numeric,
            sorting=none,
            movenames=false,    % Отключаем дублирование авторов
            doi=false,
            url=true]{biblatex}

%%% Подключение файлов bib %%%
\addbibresource{biblio/external.bib}
\addbibresource{biblio/author.bib}
\addbibresource{biblio/registered.bib}

%%% Убираем [Текст] %%%
\DeclareSourcemap{
    \maps{
        \map[overwrite]{
            \step[fieldset=media, null]
        }
    }
}

%%% Убираем квадратные скобки вокруг "и др." %%%
\renewcommand{\bibleftbracket}{}
\renewcommand{\bibrightbracket}{}

%%% Показываем всех авторов %%%
\setcounter{maxnames}{99}
\setcounter{minnames}{99}

%%% Команды для вывода списка литературы %%%
\newcommand*{\insertbibliofull}{%
    \printbibliography[keyword=bibliofull,title=Список литературы]
}
\newcommand*{\insertbiblioauthor}{%
    \printbibliography[keyword=biblioauthor,title=Публикации автора]
}
\newcommand*{\insertbiblioexternal}{%
    \printbibliography[keyword=biblioexternal,title=Список литературы]
}